\section{实验验证与有效性限制}
\label{sec:validation}

\subsection{行为一致性检验}

行为检验以标准输出中的十进制整数和换行、以及进程退出状态 0 为判据。六个输入及期望输出分别为 $-3\mapsto1$、$0\mapsto1$、$1\mapsto1$、$5\mapsto120$、$10\mapsto3628800$ 和 $20\mapsto3628800$。其中负输入先被夹取为 0，超上界输入被夹取为 10；其余输入直接进入迭代阶乘。

SysY 源、参考 LLVM IR 与参考 RV64 汇编分别执行全部六个输入，三种表示均获得 6/6 通过，共形成 18 个实现--输入组合的一致结果。C 观察形态采用相同控制流和算法，用于生成阶段观测而不构成额外语义基准。测试覆盖两个夹取分支、零次循环和多次循环路径，并同时经过外部输入输出调用。

\subsection{结论范围与外部有效性}

上述结果支持三种表示在声明输入上的可观察行为一致性，但测试集合并未穷尽整数输入空间，也未建立翻译关系的形式证明。更强的全域结论需要明确 SysY、LLVM IR 与目标机器的整数及 I/O 语义，并采用符号验证或经证明的翻译验证方法。

结构计数描述特定编译器版本、RV64 目标和静态产物，不测量动态指令数、延迟、吞吐量或能耗。优化流水线、IR 打印形式和指令选择均可能随 LLVM 版本及目标配置变化，表~\ref{tab:structure} 的数值因而属于受控环境中的观测结果。对象类型、目标属性、重定位和程序头均从实际 ELF 产物读取，以降低由文件名或驱动默认值造成的目标识别偏差。

程序执行发生于 QEMU RISC-V Linux 用户态模拟环境\cite{qemu-user-2026}，验证范围是 Linux 用户态 ABI 下的程序行为，不包含物理 RISC-V 处理器的微体系结构性能或设备特性。MLIR 与 AscendNPU IR 仅依据同行评审论文和官方文档讨论，未进入本地编译、模拟或设备实验。
